\DocumentMetadata{
  lang        = en,
%   pdfstandard = ua-2,
%   pdfstandard = a-4f, %or a-4
  pdfstandard = a-4,
  tagging=on,
  tagging-setup={math/setup=mathml-SE} 
}
\documentclass[11pt]{article}
\usepackage{paralist}
\usepackage{fullpage,graphicx}
\usepackage{bold-extra}
\usepackage{fancyhdr}
\usepackage{algorithm}
\usepackage{algpseudocode}
\usepackage[dvipsnames]{xcolor}
\input xy
\xyoption{all}
\usepackage{color,soul}
\usepackage{amsmath}
\usepackage{amssymb}
\usepackage{url}

\usepackage{listings}
\lstdefinestyle{customjava}{
  belowcaptionskip=1\baselineskip,
  breaklines=true,
  frame=L,
  xleftmargin=\parindent,
  language=java,
  showstringspaces=false,
  basicstyle=\ttfamily,
  keywordstyle=\bfseries\color{green!40!black},
  commentstyle=\itshape\color{purple!40!black},
  identifierstyle=\color{blue},
  stringstyle=\color{red!40!black},
}

\usepackage[textwidth=2cm,textsize=tiny]{todonotes}
\newcommand{\new}[1]{\textcolor{blue}{#1}}
\newcommand{\wnote}[1]{\todo[color=Yellow]{{\bf WV}: #1}}
\newcommand{\pnote}[1]{\todo[color=Lavender]{{\bf PE}: #1}}

\newcommand{\question}[1]{#1}

%ADD A % TO THE NEXT LINE, AND REMOVE THE % FROM THE FOLLOWING LINE IN ORDER TO
%MAKE THE ANSWERS DISAPPEAR
% \newcommand{\answer}[1]{\ \\ \noindent\textbf{Answer: }#1}
\newcommand{\answer}[1]{}

%\lfoot{PLEASE TURN OVER}
\renewcommand{\headrulewidth}{0pt}
\pagestyle{fancy}

\setlength{\parindent}{0pt}

\usepackage{tikz}
\usetikzlibrary{arrows,automata,shapes}
\usetikzlibrary{calc}
\usetikzlibrary{matrix}
\usetikzlibrary{positioning}
\usepackage{tkz-graph}
\newcommand{\mi}[1]{\mathit{#1}}
\newcommand{\mr}[1]{\mathrm{#1}}
\newcommand{\mt}[1]{\mathtt{#1}}
\newcommand{\mc}[1]{\mathcal{#1}}
\newcommand{\mb}[1]{\mathbb{#1}}
\newcommand{\tb}[1]{\textbf{#1}}
\newcommand{\la}{\langle}
\newcommand{\ra}{\rangle}
\newcommand{\cob}[1]{{\color{RoyalBlue} #1}}
\newcommand{\cor}[1]{{\color{RedOrange} #1}}
\newcommand{\cop}[1]{{\color{Purple} #1}}
\newcommand{\cog}[1]{{\color{ForestGreen} #1}}

\newcommand{\deft}{\mc{D}}
\newcommand{\coop}{\mc{C}}

\begin{document}

\centerline{\textsc{University of Aberdeen}} 
\centerline{Symbolic AI (CS502K)}
\centerline{Tutorial I  -- Introduction to A.I. \& Intelligent Agents} 

\hrulefill

%\textit{Please hand in your answer to item 5 below; its mark will contribute 10\% of your overall result for the course.}

\begin{enumerate}
\item Read chapter 1 of Russell and Norvig's \emph{Artificial Intelligence: A modern approach} and answer the items below in connection with the \emph{Introduction} lecture.
\begin{enumerate}
% \item 
%   \question{Together with a classmate, define the following terms:
%   \begin{enumerate}
%   \item \emph{Intelligence}
%   \item \emph{Artificial intelligence}
%   \item \emph{Agent}
%   \item \emph{Rationality}
%   \item \emph{Logical reasoning}
%   \end{enumerate}
%   }
%   \answer{
%   \begin{enumerate}
%   \item \emph{Intelligence} -- Dictionary definitions of intelligence talk about ``the capacity to acquire and apply knowledge'' or ``the faculty of thought and reason'' or ``the ability to comprehend and profit from experience.'' These are all reasonable answers, but if we want something quantifiable we would use something like ``the ability to apply knowledge in order to perform better in an environment.''
%   \item \emph{Artificial intelligence} -- We define artificial intelligence as the study and construction of agent programs that perform well in a given environment, for a given agent architecture.
%   \item \emph{Agent} -- We define an agent as an entity that takes action in response to percepts from an environment.
%   \item \emph{Rationality} -- We define rationality as the property of a system which does the ``right thing'' given what it knows. See Section 2.2 for a more complete discussion. Both describe perfect
% rationality, however; see Section 27.3.
%   \item \emph{Logical reasoning} --  we define logical reasoning as the a process of deriving new sentences from old, such that the new sentences are necessarily true if the old ones are true. (Notice that does
% not refer to any specific syntax or formal language, but it does require a well-defined
% notion of truth.)
%   \end{enumerate}
%   }

% \item 
%   \question{Turing in his original paper\footnote{Available at \url{http://phil415.pbworks.com/f/TuringComputing.pdf}} discusses several objections to his proposed enterprise and his test for intelligence. Answer:
% \begin{enumerate}
% \item Which objections still carry weight? 
% \item Are his refutations valid? 
% \item Can you think of new objections arising from developments since he wrote the paper?
% \item In the paper, he predicts that, by the year 2000, a computer will have a 30\% chance of passing a five-minute Turing Test with an unskilled interrogator. What chance do you think a computer would have today? In another 50 years?
% \end{enumerate}
%   }
%   \answer{The probability of fooling an interrogator depends on just how unskilled the interrogator
% is. One entrant in the 2002 Loebner prize competition\footnote{\url{http://www.loebner.net/Prizef/loebner-prize.html}} (which is not quite a real Turing Test) did fool one judge, although if you look at the transcript, it is hard to imagine what that judge was thinking. There certainly have been examples of a chatbot or other online agent fooling humans. For example, see See Lenny Foner’s account of the Julia chatbot at \url{foner.www.media.mit.edu/people/foner/Julia/}. We’d say the chance today is something
% like 10\%, with the variation depending more on the skill of the interrogator rather than the
% program. In 50 years, we expect that the entertainment industry (movies, video games, commercials)
% will have made sufficient investments in artificial actors to create very credible
% impersonators.}

\item 
  \question{To what extent do the following computer systems embed artificial intelligence:
  \begin{enumerate}
  \item Supermarket bar code scanners
%  \item Web search engines
%  \item Voice-activated telephone menus
  \item Internet routing algorithms that respond dynamically to the state of the network.
  \end{enumerate}}
  \answer{\begin{enumerate}
  \item Although bar code scanning is in a sense computer vision, these are not AI systems. The problem of reading a bar code is an extremely limited and artificial form of visual interpretation, and it has been carefully designed to be as simple as possible, given the hardware.
%  \item In many respects. The problem of determining the relevance of a web page to a query is a problem in natural language understanding, and the techniques are related to those presented in Chapters 22 and 23. Search engines like Ask.com, which group the retrieved pages into categories, use clustering techniques analogous to those we discuss in Chapter 20. Likewise, other functionalities provided by a search engines use intelligent techniques; for instance, the spelling corrector uses a form of data mining based on observing users’ corrections of their own spelling errors. On the other hand, the problem of indexing billions of web pages in a way that allows retrieval in seconds is a problem in database design, not in artificial intelligence.
%  \item To a limited extent. Such menus tends to use vocabularies which are very limited -- e.g. the digits, ``Yes'', and ``No'' -- and within the designers’ control, which greatly simplifies the problem. On the other hand, the programs must deal with an uncontrolled space of all kinds of voices and accents. The voice activated directory assistance programs used by telephone companies, which must deal with a large and changing vocabulary are certainly AI programs.
  \item This is borderline. There is something to be said for viewing these as intelligent agents
working in cyberspace. The task is sophisticated, the information available is partial, the
techniques are heuristic (not guaranteed optimal), and the state of the world is dynamic.
All of these are characteristic of intelligent activities. On the other hand, the task is very
far from those normally carried out in human cognition.
  \end{enumerate}
}

\item 
\question{Why would evolution tend to result in systems that act rationally? What goals are such
systems designed to achieve?}
\answer{Evolution tends to perpetuate organisms (and combinations and mutations of organisms)
that are successful enough to reproduce. That is, evolution favors organisms that can
optimize their performance measure to at least survive to the age of sexual maturity, and then
be able to win a mate. Rationality just means optimizing performance measure, so this is in
line with evolution.}

\item 
\question{Is AI a science, or is it engineering? Or neither or both? Justify your answer.}
\answer{This question is intended to be about the essential nature of the AI problem and what is
required to solve it, but could also be interpreted as a sociological question about the current
practice of AI research.
A science is a field of study that leads to the acquisition of empirical knowledge by the
scientific method, which involves falsifiable hypotheses about what is. A pure engineering
field can be thought of as taking a fixed base of empirical knowledge and using it to solve
problems of interest to society. Of course, engineers do bits of science -- e.g., they measure the
properties of building materials -- and scientists do bits of engineering to create new devices
and so on. As described in Section 1.1, the ``human'' side of AI is clearly an empirical science --
called cognitive science these days -- because it involves psychological experiments designed
out to find out how human cognition actually works. What about the the ``rational'' side?
If we view it as studying the abstract relationship among an arbitrary task environment, a
computing device, and the program for that computing device that yields the best performance
in the task environment, then the rational side of AI is really mathematics and engineering;
it does not require any empirical knowledge about the actual world -- and the actual task
environment -- that we inhabit; that a given program will do well in a given environment is a
theorem. (The same is true of pure decision theory.) In practice, however, we are interested
in task environments that do approximate the actual world, so even the rational side of AI
involves finding out what the actual world is like. For example, in studying rational agents
that communicate, we are interested in task environments that contain humans, so we have to find out what human language is like. In studying perception, we tend to focus on sensors
such as cameras that extract useful information from the actual world. (In a world without
light, cameras wouldn’t be much use.) Moreover, to design vision algorithms that are good
at extracting information from camera images, we need to understand the actual world that
generates those images. Obtaining the required understanding of scene characteristics, object
types, surface markings, and so on is a quite different kind of science from ordinary physics,
chemistry, biology, and so on, but it is still science.

In summary, AI is definitely engineering but it would not be especially useful to us if it
were not also an empirical science concerned with those aspects of the real world that affect
the design of intelligent systems for that world.}

\item 
\question{``Surely computers cannot be intelligent -- they can do only what their programmers
tell them.'' Is the latter statement true, and does it imply the former?}
\answer{This depends on your definition of ``intelligent'' and ``tell.'' In one sense computers only
do what the programmers command them to do, but in another sense what the programmers
consciously tells the computer to do often has very little to do with what the computer actually
does. Anyone who has written a program with an ornery bug knows this, as does anyone
who has written a successful machine learning program. So in one sense Samuel ``told'' the
computer ``learn to play checkers better than I do, and then play that way,'' but in another
sense he told the computer ``follow this learning algorithm'' and it learned to play. So we’re
left in the situation where you may or may not consider learning to play checkers to be s sign
of intelligence (or you may think that learning to play in the right way requires intelligence,
but not in this way), and you may think the intelligence resides in the programmer or in the
computer.}

% \item 
% \question{``Surely animals cannot be intelligent -- they can do only what their genes tell them.''
% Is the latter statement true, and does it imply the former?}
% \answer{The point of this exercise is to notice the parallel with the previous one. Whatever
% you decided about whether computers could be intelligent in the previous item, you are committed to
% making the same conclusion about animals (including humans), unless your reasons for deciding
% whether something is intelligent take into account the mechanism (programming via
% genes versus programming via a human programmer). Note that Searle makes this appeal to
% mechanism in his Chinese Room argument (see Chapter 26).}

% \item 
% \question{``Surely animals, humans, and computers cannot be intelligent—they can do only what
% their constituent atoms are told to do by the laws of physics.'' Is the latter statement true, and
% does it imply the former? \\}
% \answer{Again, the choice you make in item f above drives your answer to this question.}

\end{enumerate}

%% First-order Logic %%

\item Read chapter 2 of Russell and Norvig's \emph{Artificial Intelligence: A modern approach} and answer the items below in connection with the \emph{Intelligent Agents} lecture.

\begin{enumerate}

 \item
  \question{For each of the following assertions, say whether it is true or false and support your
answer with examples or counterexamples where appropriate:
\begin{enumerate}
% \item An agent that senses only partial information about the state cannot be perfectly rational.
% \item There exist task environments in which no pure reflex agent can behave rationally.
% \item  There exists a task environment in which every agent is rational.
% \item  The input to an agent program is the same as the input to the agent function.
% \item  Every agent function is implementable by some program/machine combination.
%\item  Suppose an agent selects its action uniformly at random from the set of possible actions.
%There exists a deterministic task environment in which this agent is rational.
%\item  It is possible for a given agent to be perfectly rational in two distinct task environments.
\item  Every agent is rational in an unobservable environment.
\item  A perfectly rational poker-playing agent never loses.
\end{enumerate}}
  \answer{
  \begin{enumerate}
% \item False. Perfect rationality refers to the ability to make good decisions given the sensor
% information received.
% \item True. A pure reflex agent ignores previous percepts, so cannot obtain an optimal state
% estimate in a partially observable environment. For example, correspondence chess is
% played by sending moves; if the other player’s move is the current percept, a reflex agent
% could not keep track of the board state and would have to respond to, say, ``a4'' in the
% same way regardless of the position in which it was played.
% \item True. For example, in an environment with a single state, such that all actions have the
% same reward, it doesn’t matter which action is taken. More generally, any environment
% that is reward-invariant under permutation of the actions will satisfy this property.
% \item False. The agent function, notionally speaking, takes as input the entire percept sequence
% up to that point, whereas the agent program takes the current percept only.
% \item False. For example, the environment may contain Turing machines and input tapes and
% the agent’s job is to solve the halting problem; there is an agent function that specifies
% the right answers, but no agent program can implement it. Another example would be
% an agent function that requires solving intractable problem instances of arbitrary size in
% constant time.
%\item True. This is a special case of (iii); if it doesn’t matter which action you take, selecting
%randomly is rational.
% \item True. For example, we can arbitrarily modify the parts of the environment that are
% unreachable by any optimal policy as long as they stay unreachable.
\item False. Some actions are stupid -- and the agent may know this if it has a model of the
environment -- even if one cannot perceive the environment state.
\item False. Unless it draws the perfect hand, the agent can always lose if an opponent has
better cards. This can happen for game after game. The correct statement is that the
agent’s expected winnings are nonnegative.
  \end{enumerate}}

\item
\question{Give a PEAS description of the task environments below and characterize them in terms of the properties discussed in section 2.3.2 of the textbook
\begin{enumerate}
\item Shopping for used AI books on the Internet
\item Bidding on an item at an auction.
\end{enumerate}
}
\answer{\begin{enumerate}
\item Partially observable, deterministic, sequential, static, discrete, single agent. This can be
multi-agent and dynamic if we buy books via auction, or dynamic if we purchase on a
long enough scale that book offers change.
\item Fully observable, strategic, sequential, static, discrete, multi-agent.
\end{enumerate}
}

\item
\question{This exercise explores the differences between agent functions and agent programs.
\begin{enumerate}
%\item Can there be more than one agent program that implements a given agent function? Give an example, or show why one is not possible.
\item Are there agent functions that cannot be implemented by any agent program?
% \item Given a fixed machine architecture, does each agent program implement exactly one
% agent function?
% \item Given an architecture with $n$ bits of storage, how many different possible agent programs
% are there?
\item Suppose we keep the agent program fixed but speed up the machine by a factor of two.
Does that change the agent function?
\end{enumerate}
}
\answer{
\begin{enumerate}
%\item Yes; take any agent program and insert null statements that do not affect the output.
\item Yes; the agent function might specify that the agent print true when the percept is a
Turing machine program that halts, and false otherwise. (Note: in dynamic environments,
for machines of less than infinite speed, the rational agent function may not be
implementable; e.g., the agent function that always plays a winning move, if any, in a
game of chess.)
% \item Yes; the agent’s behavior is fixed by the architecture and program.
% \item There are $2^n$ agent programs, although many of these will not run at all. (Note: Any
% given program can devote at most $n$ bits to storage, so its internal state can distinguish
% among only $2^n$ past histories. Because the agent function specifies actions based on percept
% histories, there will be many agent functions that cannot be implemented because
% of lack of memory in the machine.)
\item It depends on the program and the environment. If the environment is dynamic, speeding
up the machine may mean choosing different (perhaps better) actions and/or acting
sooner. If the environment is static and the program pays no attention to the passage of
elapsed time, the agent function is unchanged.
\end{enumerate}
}

\item 
\question{Discuss possible agent programs for each of the following stochastic environments:
\begin{enumerate}
\item Murphy's law: twenty-five percent of the time, the Suck action fails to clean the floor if
it is dirty and deposits dirt onto the floor if the floor is clean. How is your agent program
affected if the dirt sensor gives the wrong answer 10\% of the time?
\item Small children: At each time step, each clean square has a 10\% chance of becoming
dirty. Can you come up with a rational agent design for this case?
\end{enumerate}
}
\answer{
\begin{enumerate}
\item For a reflex agent, this presents no additional challenge, because the agent will continue
to $\mathit{Suck}$ as long as the current location remains dirty. For an agent that constructs a
sequential plan, every $\mathit{Suck}$ action would need to be replaced by ``$\mathit{Suck}$ until clean.''
If the dirt sensor can be wrong on each step, then the agent might want to wait for a
few steps to get a more reliable measurement before deciding whether to $\mathit{Suck}$ or move
on to a new square. Obviously, there is a trade-off because waiting too long means
that dirt remains on the floor (incurring a penalty), but acting immediately risks either
dirtying a clean square or ignoring a dirty square (if the sensor is wrong). A rational
agent must also continue touring and checking the squares in case it missed one on a
previous tour (because of bad sensor readings). it is not immediately obvious how the
waiting time at each square should change with each new tour. These issues can be
clarified by experimentation, which may suggest a general trend that can be verified
mathematically. This problem is a partially observable Markov decision process -- see
Chapter 17. Such problems are hard in general, but some special cases may yield to
careful analysis.
\item In this case, the agent must keep touring the squares indefinitely. The probability that
a square is dirty increases monotonically with the time since it was last cleaned, so the
rational strategy is, roughly speaking, to repeatedly execute the shortest possible tour of
all squares. (We say ``roughly speaking'' because there are complications caused by the
fact that the shortest tour may visit some squares twice, depending on the geography.)
This problem is also a partially observable Markov decision process.
\end{enumerate}
}
\end{enumerate}

% \item (\textbf{Hand-in Exercise}) This part of the tutorial will be available on MyAberdeen, as part of the course's summative assessment. You have been tasked with designing a simple reflex software agent to control the heating and ventilation of a room. The can control the setting of the heating which can be maximum, moderate or no heating, and the setting of the ventilation which can be strong, weak or no ventilation. The agent has access to the temperature outside the room (i.e., low, medium, high), to the level of humidity in the room is (i.e., good or dry) and the price band of energy for both heating and ventilation (i.e., expensive, moderate, cheap). You should
% \begin{enumerate}
%     \item List the sensors the agent has access to.\hfill [1 mark]
%     \item List the actuators the agent has access to.\hfill [1 mark]
%     \item Provide at least 3 condition-action rules the agent should use to decide on which action to take. Your rules should use all sensors and actuators from above\hfill [3 marks]
% \end{enumerate}
% \answer{\begin{enumerate}
%     \item The sensors detect temperature outside, humidity in the room and current price band of gas.
%     \item The actuators control the amount of heating and ventilation.
%     \item Some sample rules are: \\
%      \begin{tabular}{ll}
%          if & temperature is low and humidity is dry and price band is cheap and \\
%             & heating is off and ventilation is no \\
%          then & set heating to maximum and ventilation to strong \\
%          if & temperature is high and humidity is dry and price band is cheap and \\
%             & heating is on and ventilation is no \\
%          then & set heating to no-heating and ventilation to strong \\
%          if & temperature is high and humidity is dry and price band is expensive and \\
%             & heating is medium and ventilation is no \\
%          then & set heating to no-heating and ventilation to weak
%     \end{tabular}
% \end{enumerate}
% \textbf{Marking scheme}
% \begin{enumerate}
%     \item 0.4 mark per first sensor correctly identified; 0.3 marks per subsequent sensors correctly identified.
%     \item 0.5 mark per actuator correctly identified.
%     \item 1 mark per correct rule which should list all sensors and actuators, although no marks deducted if the status of actuators is missing from the ``if-part'' of the rules. Rules, however, must ``make sense'' in that they should describe realistic situations and appropriate actions for these situations.
% \end{enumerate}}

\item 
\question{Familiarise yourself with the code base for the algorithms presented in the AIMA textbook; the code is available at \url{https://github.com/aimacode/aima-python}. You have two options to work with aima-python: (1) install Anaconda + Jupyter on your computer--recommended; (2) run aima-python directly online in a Binder Environment – worth trying.

\begin{enumerate}
\item Start exploring aima-python with the intro Jupyter notebook (e.g. \url{https://github.com/aimacode/aima-python/blob/master/intro.ipynb}). If you have a local aima-python (as recommended above) you start with your local version of the intro Jupyter notebook.
\item The intro notebook has a link to another notebook called index (e.g. \url{https://github.com/aimacode/aima-python/blob/master/index.ipynb}). If you have a local aima-python (as recommended above) you start with your local version of the index Jupyter notebook.
\end{enumerate}
}
\answer{Follow the instructions in the aima-python code repository.}

\item 
\question{The agents Jupyter notebook (e.g. \url{https://github.com/aimacode/aima-python/blob/master/agents.ipynb}) has a few examples of agents and environments – i.e. a BlindDog in a park. You will notice that the agents.ipynb imports code from agents.py. This is a common feature in aima-python where Jupyter notebooks use python libraries that actually implement the algorithms from the AIMA textbook.

\begin{enumerate}
\item Inspect the code for the abstract Agent and Environment classes which are defined in the \texttt{agents.py} library.
\item Starting with a simple BlindDog agent and a Park environment, the agents.ipynb notebook has multiple implementations of BlindDog agents and their Park environments increasing the complexity of the environment. You will learn how complex environments increases the sophistication of the agents e.g. EnergeticBlindDog.
\item The agents.ipynb also shows an example of the Wumpus environment. We will only learn building functionality to reason about individual actions an agent in the Wumpus world can take at any given stage of the agent's journey in the Wumpus world. A fully functioning agent for the Wumpus environment is way more complex than building a BlindDog. Why?
\end{enumerate}
}

\answer{
Follow the instructions in the aima-python code repository.
}

\end{enumerate} % end of top-level enumeration

\end{document}






























